
\documentclass[final,t]{beamer}

\usepackage[size=a0,orientation=portrait,scale=0.98]{beamerposter}
\usepackage[utf8]{inputenc}
\usepackage[T1]{fontenc}
\usepackage{lmodern}
\usepackage[spanish,es-tabla]{babel}
\usepackage[table]{xcolor}
\usepackage{graphicx}
\usepackage{amsmath}
\usepackage{booktabs}
\usepackage{tabularx}
\usepackage{array}
\usepackage{tikz}
\usetikzlibrary{positioning,calc}
\graphicspath{{docs/}{../}}

% Overleaf: subir tambien taxonomia-final-merged-after-rework-treemap-poster-wide.pdf en la raiz del proyecto o en docs/.
\setbeamertemplate{navigation symbols}{}
\setbeamertemplate{headline}{}
\setbeamertemplate{footline}{}
\setbeamersize{text margin left=14mm,text margin right=14mm}
\renewcommand{\familydefault}{\sfdefault}
\setlength{\parindent}{0pt}
\setlength{\parskip}{0.20em}
\renewcommand{\arraystretch}{1.12}
\newcolumntype{Y}{>{\raggedright\arraybackslash}X}

\definecolor{paperText}{HTML}{1A1A1A}
\definecolor{paperMuted}{HTML}{5F6368}
\definecolor{paperRule}{HTML}{B7B7B7}
\definecolor{paperAccent}{HTML}{1F4E79}
\definecolor{paperAccentLight}{HTML}{EAF2F8}
\definecolor{paperSoft}{HTML}{F7F7F7}
\definecolor{barOne}{HTML}{315C8C}
\definecolor{barTwo}{HTML}{5F84B2}
\definecolor{barThree}{HTML}{8EABCC}
\definecolor{barFour}{HTML}{B5C7DC}
\definecolor{barOther}{HTML}{D9E2EC}

\setbeamercolor{background canvas}{bg=white}
\setbeamercolor{normal text}{fg=paperText,bg=white}
\setbeamerfont{normal text}{size=\large}

\newcommand{\articleblock}[2]{%
  {\color{paperAccent}\LARGE\bfseries #1}\par
  \vspace{0.05cm}{\color{paperRule}\hrule height 1.0pt}\vspace{0.08cm}
  {#2}\par\vspace{0.08cm}
}
\newcommand{\smallnote}[1]{{\color{paperMuted}\normalsize #1}}
\newcommand{\compactnote}[1]{{\color{paperMuted}\small #1}}

\newcommand{\metriccell}[3]{%
  \begin{minipage}[t]{0.31\linewidth}
    \centering
    {\fontsize{58}{62}\selectfont\bfseries #1}\par
    {\Large\bfseries #2}\par
    \smallnote{#3}
  \end{minipage}%
}

\newcommand{\barrowwide}[5]{%
  \noindent\begin{minipage}[t]{0.28\linewidth}\raggedright\Large #1\end{minipage}%
  \hfill
  \begin{minipage}[t]{0.58\linewidth}
    \vspace{0.20em}%
    \begin{tikzpicture}[x=1cm,y=1cm]
      \pgfmathsetmacro{\barwidth}{#3*(#2)/104.0}%
      \draw[fill=#5, draw=none] (0,0) rectangle (\barwidth,0.52);
      \draw[draw=paperRule, line width=0.35pt] (0,0) rectangle (#3,0.52);
    \end{tikzpicture}%
  \end{minipage}%
  \hfill
  \begin{minipage}[t]{0.10\linewidth}\raggedleft\Large\bfseries #4\end{minipage}\par\vspace{0.16cm}%
}

\newcommand{\funnelrow}[5]{%
  \noindent\begin{minipage}[t]{0.30\linewidth}\raggedright\Large #1\end{minipage}%
  \hfill
  \begin{minipage}[t]{0.56\linewidth}
    \vspace{0.20em}%
    \begin{tikzpicture}[x=1cm,y=1cm]
      \pgfmathsetmacro{\barwidth}{#3*(#2)/104.0}%
      \draw[fill=#5, draw=none] (0,0) rectangle (\barwidth,0.50);
      \draw[draw=paperRule, line width=0.35pt] (0,0) rectangle (#3,0.50);
    \end{tikzpicture}%
  \end{minipage}%
  \hfill
  \begin{minipage}[t]{0.11\linewidth}\raggedleft\Large\bfseries #4\end{minipage}\par\vspace{0.14cm}%
}

\begin{document}
\begin{frame}[t]

\vspace*{0.08cm}
\begin{center}
  
\includegraphics[height=2.8cm]{presentacion_poster/ufro-logo.png}\par
  \vspace{0.14cm}
  {\fontsize{46}{52}\selectfont\bfseries Taxonomía de motivos de retrabajo en pull requests de agentes de IA}\par
  \vspace{0.14cm}
  {\Large Javier Alcalde y Diego Labrín \quad | \quad Profesor: Dr. Pablo Valenzuela}\par
  \vspace{0.06cm}
  {\large Departamento de Computación e Informática, Universidad de La Frontera \quad | \quad Minería de Repositorios de Software}\par
  \vspace{0.10cm}
  {\color{paperRule}\rule{0.94\textwidth}{1.5pt}}
\end{center}

\vspace{0.12cm}
\begin{center}
\begin{minipage}{0.94\textwidth}

\articleblock{Pregunta central y objetivos específicos}{
Los agentes de IA pueden producir PRs mergeables, pero la integración suele requerir retrabajo humano previo. En AIDev analizamos PRs \textbf{mergeados después de retrabajo observable} para identificar los motivos que aparecen antes del merge.

\vspace{0.18cm}
\begin{tabularx}{\linewidth}{>{\bfseries}p{6.1cm}Y}
Pregunta central & ¿Qué motivos de retrabajo humano emergen en pull requests de agentes de IA antes de su integración? \\
Objetivo 1 & Definir una población operacional auditable a partir del dataset AIDev, manteniendo trazabilidad entre población, muestra, tarjetas y categorías preliminares. \\
Objetivo 2 & Implementar una muestra estratificada reproducible por agente de IA. \\
Objetivo 3 & Generar tarjetas con evidencia textual humana y aplicar una primera etapa de card sorting manual coherente, con cita textual y justificación breve. \\
\end{tabularx}
}

\articleblock{Diseño metodológico y trazabilidad}{
\begin{center}
\begin{tikzpicture}[x=1cm,y=1cm]
  \node[draw=paperAccent, fill=paperAccentLight, rounded corners=5pt, minimum width=14.6cm, minimum height=2.0cm, align=center] (prep) at (0,0) {\bfseries Preparation\\\smallnote{filtros, muestra estratificada, tarjetas}};
  \node[draw=paperAccent, fill=white, rounded corners=5pt, minimum width=14.6cm, minimum height=2.0cm, align=center] (exec) at (17.2,0) {\bfseries Execution\\\smallnote{card sorting abierto}};
  \node[draw=paperAccent, fill=paperAccentLight, rounded corners=5pt, minimum width=14.6cm, minimum height=2.0cm, align=center] (ana) at (34.4,0) {\bfseries Analysis\\\smallnote{taxonomía, métricas y soundness}};
  \draw[-latex, line width=1.35pt, paperAccent] (prep) -- (exec);
  \draw[-latex, line width=1.35pt, paperAccent] (exec) -- (ana);
\end{tikzpicture}
\end{center}

\vspace{0.18cm}
\begin{tabularx}{\linewidth}{>{\bfseries}p{4.6cm}Y>{\bfseries}p{4.6cm}Y}
Inclusión & PR cerrado y mergeado, \texttt{merged\_at} no nulo, \texttt{commit\_count > 1}, \texttt{human\_comment\_count > 0}. & Evidencia & Comentarios, reviews y timeline cruzados con la muestra; una tarjeta por PR. \\
Exclusión & PRs abiertos, cerrados sin merge, sin comentario humano observable o rechazo definitivo. & Validez & Cada asignación debe responder la RQ y estar respaldada por evidencia textual o marcarse como insuficiente. \\
\end{tabularx}
}

\articleblock{Datos, filtros y muestra}{
\begin{center}
\metriccell{33.596}{PRs en AIDev}{universo bruto}\hfill
\metriccell{3.166}{PRs operacionales}{merged\_after\_rework}\hfill
\metriccell{300}{PRs muestreados}{seed 20260510}
\end{center}

\vspace{0.30cm}
\funnelrow{Universo AIDev}{104}{31.0}{100,00\%}{barOne}
\funnelrow{PRs mergeados}{74.3}{31.0}{71,48\%}{barTwo}
\funnelrow{Con commits adicionales}{21.3}{31.0}{20,49\%}{barThree}
\funnelrow{Población operacional}{9.8}{31.0}{9,42\%}{barFour}
\funnelrow{Muestra estratificada}{0.9}{31.0}{0,89\%}{barOther}

\vspace{0.06cm}
Asignación proporcional por agente: \(n_h = \operatorname{round}(n \times N_h / N)\), con \(n=300\), \(N=3.166\), mínimo por estrato y semilla fija. Cuotas: Copilot 145, Devin 86, OpenAI\_Codex 45, Cursor 17 y Claude\_Code 7.
}


\articleblock{Resultado principal: treemap de taxonomía final}{
Las cuatro primeras familias explican \textbf{228 de 300 casos} (\textbf{76\%}). El treemap muestra la taxonomía final: cada área representa la frecuencia de un motivo de retrabajo humano pre-merge y las leyendas explicitan el tipo de rechazo temporal observado en comentarios/reviews.

\vspace{0.10cm}
\begin{center}
% Treemap original (layout correcto, sin comas decimales).
% 1. Descarga docs/taxonomia-final-merged-after-rework-treemap.svg
% 2. Conviértelo a PDF (ej: convertio.co) y súbelo a Overleaf como .pdf
% 3. O compila con LuaLaTeX/XeLaTeX (soporte SVG nativo con paquete svg).
\includegraphics[width=0.87\linewidth]{taxonomia-final-merged-after-rework-treemap.pdf}
\end{center}
}


\articleblock{Discusión e implicancias reforzadas}{
\begin{tabularx}{\linewidth}{>{\bfseries}p{5.2cm}YY}
\toprule
Foco & Lectura del hallazgo & Implicancia práctica \\
\midrule
CI y calidad & P1 concentra 104 casos: fallos automáticos, formato, lint y cobertura. La aceptación no depende solo de generar código plausible, sino de satisfacer contratos verificables del repositorio. & Los agentes deberían ejecutar checks locales, interpretar fallos de CI, corregir estilo/pruebas y explicar qué validaciones pasaron antes de pedir revisión humana. \\
Impl. & P2 muestra retrabajo semántico: casos borde, compatibilidad, rendimiento, recursos y comportamiento frontend. El problema no es solo sintaxis, sino adecuación funcional al contexto real de uso. & La evaluación debe medir trayectoria de integración, commits adicionales, cambios solicitados y feedback humano; la tasa de merge aislada sobreestima autonomía. \\
Arq. y alcance & P3 indica que una solución técnicamente posible puede ser demasiado amplia, duplicada o inconsistente con el diseño local. El retrabajo corrige encaje arquitectónico, no solo errores. & Las herramientas IA necesitan memoria de repositorio: patrones arquitectónicos, convenciones, dependencias activas, owners, locks y límites explícitos del issue. \\
Gobernanza & P4 muestra que el merge depende de aprobación, coordinación, orden de integración y reglas sociales del proyecto. La integración es un proceso colectivo, no un resultado binario. & Los mantenedores pueden reducir retrabajo repetitivo haciendo explícitos criterios de aceptación, reglas de contribución, checks obligatorios y dependencias entre PRs. \\
\bottomrule
\end{tabularx}

\vspace{0.18cm}
\begin{tabularx}{\linewidth}{>{\bfseries}p{5.2cm}Y>{\bfseries}p{5.2cm}Y}
Lectura transversal & El retrabajo pre-merge funciona como una capa de alineación sociotécnica: transforma una contribución generada por IA en una contribución aceptable para una comunidad y un repositorio concretos. & Riesgo metodológico & Si solo contamos PRs mergeados, perdemos el costo oculto de revisión, coordinación y corrección que hizo posible ese merge. \\
Diseño de agentes & La mejora no debería limitarse a mejores modelos de generación; también requiere loops de auto-verificación, lectura de feedback, respeto por alcance y adaptación a normas locales. & Uso para equipos & La taxonomía puede guiar checklists de revisión, plantillas de PR, reglas de CI y prompts internos para prevenir las familias de retrabajo más frecuentes. \\
\end{tabularx}
}

\articleblock{Alcance, límites y trabajo futuro}{
\begin{tabularx}{\linewidth}{>{\bfseries}p{4.8cm}YYY}
\toprule
Dimensión & Estado en esta entrega & Uso previsto & Precaución \\
\midrule
Hecho & Población operacional, muestra estratificada, tarjetas con evidencia textual humana, card sorting manual y taxonomía final trazable. & Responder la pregunta central: qué motivos de retrabajo emergen antes de la integración. & La taxonomía clasifica el motivo principal; algunos PRs contienen más de una señal de retrabajo. \\
No implementado & Medición del tiempo de retrabajo hasta merge y comparación temporal por categoría. & Queda como trabajo futuro usando creación del PR, primera evidencia de retrabajo y \texttt{merged\_at}. & No se reportan conclusiones temporales en este póster; mencionarlo como resultado sería sobreinterpretar el alcance. \\
Esfuerzo observable & Comentarios/reviews, citas textuales, commits adicionales y señales de CI o configuración presentes en la evidencia. & Permite describir tipos de intervención humana y técnica, pero no aún estimar duración del retrabajo. & Más comentarios o commits no siempre equivalen a más complejidad; requieren análisis contextual. \\
\bottomrule
\end{tabularx}

\vspace{0.10cm}
\begin{tabularx}{\linewidth}{>{\bfseries}p{4.8cm}Y>{\bfseries}p{4.8cm}Y}
Límites & Una tarjeta se asigna a una subcategoría principal; la muestra caracteriza AIDev y no generaliza automáticamente. & Reprod. & Población, muestra, tarjetas, taxonomía y scripts están versionados; seed \texttt{20260510}. \\
Refs. & Liang, Yang y Myers, ICSE 2024; Vaithilingam, Zhang y Glassman, CHI 2022; Barke, James y Polikarpova, 2022. & Repo & \texttt{github.com/Diego-LC/proyecto-semestral-msr-aidev}. \\
\end{tabularx}
}

\end{minipage}
\end{center}

\end{frame}
\end{document}
